\documentclass[12pt,a4paper]{report}

% -------------------------
% Packages
% -------------------------
\usepackage[utf8]{inputenc}
\usepackage[T1]{fontenc}
\usepackage[french]{babel}

\usepackage{geometry}
\geometry{
    left=3cm,
    right=2.5cm,
    top=2.5cm,
    bottom=2.5cm
}

\usepackage{setspace}
\onehalfspacing

\usepackage{graphicx}
\usepackage{float}
\usepackage{caption}
\usepackage{subcaption}

\usepackage{amsmath, amssymb}
\usepackage{booktabs}
\usepackage{array}
\usepackage{longtable}

\usepackage{xcolor}
\usepackage{hyperref}
\usepackage{url}

\usepackage{fancyhdr}
\usepackage{titlesec}
\usepackage{enumitem}

% -------------------------
% Configuration
% -------------------------
\hypersetup{
    colorlinks=true,
    linkcolor=black,
    urlcolor=blue,
    citecolor=black,
    % pdftitle={Titre du rapport},
    % pdfauthor={Nom Prénom}
}

\pagestyle{fancy}
\fancyhf{}
\fancyhead[L]{Efficacité et compression}
\fancyhead[R]{Pier-Luc Dupéré}
\fancyfoot[C]{\thepage}

\setlength{\parindent}{0pt}
\setlength{\parskip}{6pt}


% -------------------------
% Début du document
% -------------------------
\begin{document}

% -------------------------
% Page de garde
% -------------------------
\begin{titlepage}

    \begin{center}

        \includegraphics[width=0.50\textwidth]{img/logo-universite.png}
        
        \vspace{1cm}

        {\Large \textbf{Université du Québec à Rimouski}}\\[0.3cm]
        {\large Département de mathématiques, d'informatique et de génie}\\[0.3cm]
        {\large 8INF976 - Sujet spécial en intelligence artificielle}\\[1.5cm]

        \rule{\textwidth}{0.5pt}\\[0.4cm]

        {\Huge \textbf{Efficacité et compression}}\\[0.4cm]
        {\Large Rapport de lecture}\\[0.4cm]

        \rule{\textwidth}{0.5pt}\\[2cm]

        \begin{flushleft}
            \textbf{Réalisé par :} Pier-Luc Dupéré \\[0.4cm]
            \textbf{Encadré par :} Mehdi Adda, PhD 
        \end{flushleft}

        \vfill
    \end{center}

\end{titlepage}

\tableofcontents
% \listoffigures
% \listoftables

% -------------------------
% Corps du rapport
% -------------------------
\clearpage
\pagenumbering{arabic}

\chapter{Introduction}

Présentez ici le contexte général du sujet, la problématique et l'intérêt du travail.

\section{Contexte}
Décrivez le contexte universitaire, professionnel ou scientifique du rapport.

\section{Problématique}
Présentez la question principale à laquelle votre travail tente de répondre.

\section{Objectifs}
Les objectifs de ce rapport sont les suivants :

\begin{itemize}
    \item présenter le contexte du projet ;
    \item expliquer la méthodologie utilisée ;
    \item analyser les résultats obtenus ;
    \item proposer des perspectives d'amélioration.
\end{itemize}

\section{Organisation du rapport}
Ce rapport est organisé en plusieurs chapitres. Le premier chapitre présente le contexte, le deuxième décrit la méthodologie et le troisième analyse les résultats.

\chapter{État de l'art}

\section{Présentation générale}
Présentez les concepts et les travaux existants liés à votre sujet.

\section{Travaux connexes}
Discutez des recherches, méthodes ou solutions déjà proposées.

\section{Synthèse}
Comparez les différentes approches et expliquez le choix retenu.

\chapter{Méthodologie}

\section{Matériel et outils}
Présentez les logiciels, outils, équipements ou données utilisés.

\begin{table}[H]
    \centering
    \caption{Outils utilisés}
    \begin{tabular}{lll}
        \toprule
        Outil & Version & Utilisation \\
        \midrule
        LaTeX & 2025 & Rédaction du rapport \\
        Python & 3.x & Analyse des données \\
        Git & 2.x & Gestion du projet \\
        \bottomrule
    \end{tabular}
\end{table}

\section{Méthode de travail}
Expliquez les différentes étapes suivies pour réaliser le travail.

\section{Architecture ou processus}
Insérez ici un schéma, une architecture ou un diagramme.

\begin{figure}[H]
    \centering
    \includegraphics[width=0.7\textwidth]{img/schema.png}
    \caption{Schéma général du projet}
    \label{fig:schema}
\end{figure}

La figure~\ref{fig:schema} présente l'organisation générale du projet.

\chapter{Réalisation}

\section{Développement}
Décrivez les étapes de réalisation du projet.

\section{Implémentation}
Vous pouvez insérer du code ou des équations.

Exemple d'équation :

\begin{equation}
    E = mc^2
    \label{eq:energie}
\end{equation}

\section{Difficultés rencontrées}
Présentez les problèmes rencontrés ainsi que les solutions adoptées.

\chapter{Résultats et discussion}

\section{Résultats obtenus}
Présentez et analysez vos résultats.

\begin{figure}[H]
    \centering
    \includegraphics[width=0.75\textwidth]{img/resultats.png}
    \caption{Résultats obtenus}
    \label{fig:resultats}
\end{figure}

\section{Interprétation}
Expliquez la signification des résultats et comparez-les avec les objectifs initiaux.

\section{Limites du travail}
Présentez les limites de votre étude ou de votre projet.

\chapter{Conclusion et perspectives}

Ce rapport a permis de présenter le contexte, la méthodologie et les résultats du travail réalisé.

Les principaux résultats montrent que les objectifs fixés ont été atteints. Cependant, certaines améliorations restent possibles, notamment :

\begin{itemize}
    \item améliorer la précision des résultats ;
    \item tester la solution sur un ensemble de données plus important ;
    \item ajouter de nouvelles fonctionnalités ;
    \item optimiser les performances.
\end{itemize}

% -------------------------
% Bibliographie
% -------------------------
% \begin{thebibliography}{9}

% \bibitem{latex}
% Leslie Lamport,
% \textit{LaTeX: A Document Preparation System},
% Addison-Wesley, 1994.

% \bibitem{article}
% Nom de l'auteur,
% ``Titre de l'article'',
% \textit{Nom de la revue}, volume, numéro, année.

% \bibitem{site}
% Nom de l'auteur ou de l'organisation,
% \textit{Titre de la page web},
% \url{https://www.exemple.com},
% consulté le 18 septembre 2026.

% \end{thebibliography}

% -------------------------
% Annexes
% -------------------------
% \appendix

% \chapter{Annexe A : Code source}

% Présentez ici des extraits de code, des tableaux supplémentaires ou des documents complémentaires.

% \chapter{Annexe B : Informations supplémentaires}

% Ajoutez ici les éléments qui ne sont pas indispensables à la compréhension du rapport principal.

\end{document}
